% Pursuing Stacks (cotes 134-2 to 134-8): archive cote 134-4, site page 62 (PDF page 63), typescript page 150.
% typescript. Draft compared: ps3.tex:3153-3218.
% AI-transcribed by Claude Opus 5.5 (claude-opus-5-5, adaptive thinking, effort high), Message Batches API, batch ps2 (msgbatch_01NVwr1Yuo8qJf5rykeGBXPK), 2026-11-12; not yet checked by eye.
\scanpage{134-4/62}
\mspage{150}

A topos \hand{$\mathcal{A}$} as in the proposition (namely strictly totally aspheric)
need not be \deleted{a} modeliz\over{er}{ing}, i.e. \xout{xxxxxxxx} the Lawvere element L need not
be aspheric, or what amounts to the same because of $h_W = h_M$ and L being
\margin{x} a \uline{contrac}tor, L need not be connected: take \hand{$\mathcal{A}$} = (Sets) ! I suspect
\margin{X} though this to be the only counterexample (up to equivalence)\hand{(*)}.For \hand{$\mathcal{A}$} to
be a modelizer, we need only find an object in A which has got two
\uline{distinct} sections (because then \unsure{they|thay} must be disjoint i.e. $e_0 \wedge e_1 = \emptyset_{\hand{\mathcal{A}}}$ ,
because e has only the full and the ``empty'' subobject, as a consequence
of every ``non empty'' object of \hand{$\mathcal{A}$} having a section), thus getting a
``homotopy interval'' (more specifically, a \inserted[typed]{separated and} relatively aspheric \xout{and} \unsure{A} one)
as requested for \hand{$\mathcal{A}$} to be a modelizer. Now for any \inserted[typed]{would-be} test category met
with so far (except precisely \hand{$\Delta$}$^f$ and the like, which are \xout{xxxx} \uline{not}
test categories but only weak ones), this condition that there are objects
in A which have more than just one``point''\xout{xx} \hand{(}\xout{i.e.}= section\hand{)}, is trivially
verified.

In case of a topos of the type A\^{}, the notion of weak equivalence in
A\^{} can be described (independently of cohomological notions) in terms of
the notion of weak equivalence in (Cat), more precisely
\[ W_A \overset{\text{def}}{=} W_{A\hat{}} = i_A^{-1}(\underline{W}) , \]
where
\[ i_A :\ A\hat{} \hand{\to} (\mathrm{Cat}) , \quad F \hand{\mapsto} A_{/F} \]
is the canonical functor, and where
\[ \underline{W} = W_{(\mathrm{Cat})} \]
is the set of weak equivalences in (Cat). These of \unsure{course|dourse}, for the time
being, are defined in terms of \xout{(non commutative)} cohomology (including
a bit of non commutative one in dimension\xout{s} 1 ...). We may however start
with any $\underline{W} \subset Fl((\mathrm{Cat}))$ and look at which formal properties on $\underline{W}$ (satis-%
fied for usual weak equivalences) allow our arguments to go through, in
various circumstances. We may make a list of those which have been used
today, and go on this way a little longer, with the expectation we'll
finally wind up with an axiomatic characterization of weak equivalences
i.e. of $\underline{W}$,in terms of the category (Cat)\unsure{,} say.
a) (Pour mémoire !) $\underline{W}$ is saturated (cf. page 101).
b) $\underline{W}$ is 0-admissible , i.e. if f: C $\hand{\to}$ C' is in $\underline{W}$ , $\hand{\pi}_0$(f) is bijective

c) $\hand{\Delta}_1$ is $\underline{W}$-aspheric over e = $\hand{\Delta}_0$ , i.e. for every C in (Cat), the
projection C x$\hand{\Delta}_1 \hand{\longrightarrow}$ C is in $\underline{W}$ .
d) Any C in (Cat) which has a final element is $\underline{W}$-aspheric i.e.
C $\hand{\to}$ e is in $\underline{W}$ .

\margin{62}
